\subsection{Applications: I. Vector spaces of continuous real functions}

Let X be a nonempty compact space, H a linear subspace of the Banach space \(\mathcal{C}(\mathrm{X};\mathbf{R})\) that contains the constants and separates the points of X (GT, X, §4, No.~1, Def. 1). We equip \(\mathcal{H}\) with the normed space topology induced by that of \(\mathcal{C}(\mathrm{X};\mathbf{R})\) , and denote by \(\mathcal{H}'\) the dual of this normed space. For every \(x\in \mathrm{X}\) , the mapping \(f\mapsto f(x)\) is a continuous linear form on \(\mathcal{H}\) (the restriction to \(\mathcal{H}\) of the Dirac measure \(\varepsilon_{x}\) ), thus is an element of \(\mathcal{H}'\) that will be denoted \(i_{\mathcal{H}}(x)\) , so that

\[
\langle f, i _ {\mathcal {H}} (x) \rangle = f (x)\tag{9}
\]

for every function \(f \in \mathcal{H}\) and every \(x \in X\) .

The mapping \(i_{\mathcal{H}}\) of X into \(\mathcal{H}'\) is injective and continuous when \(\mathcal{H}'\) is equipped with the weak topology \(\sigma(\mathcal{H}', \mathcal{H})\) ; the second assertion follows at once from the definitions and (9); as for the first, note that if \(x, x'\) are two distinct points of X, by hypothesis there exists a function \(h \in \mathcal{H}\) such that \(h(x) \neq h(x')\) , therefore, by (9), \(\langle h, i_{\mathcal{H}}(x) \rangle \neq \langle h, i_{\mathcal{H}}(x') \rangle\) and a fortiori \(i_{\mathcal{H}}(x) \neq i_{\mathcal{H}}(x')\) . The image \(i_{\mathcal{H}}(\mathrm{X})\) is therefore a compact subset of \(\mathcal{H}'\) (for the weak topology), and \(i_{\mathcal{H}}\) is a homeomorphism of X onto \(i_{\mathcal{H}}(\mathrm{X})\) .

PROPOSITION 4. --- (i) The closed convex envelope C of \(i_{\mathcal{H}}(X)\) in \(\mathcal{H}'\) (for the weak topology \(\sigma(\mathcal{H}', \mathcal{H})\) ) is compact.

\begin{enumerate}
\def\labelenumi{(\roman{enumi})}
\setcounter{enumi}{1}
\tightlist
\item
  For a point \(i_{\mathcal{H}}(x)\) to be an extremal point of \(C\) , it is necessary and sufficient that the only positive measure \(\lambda\) on \(X\) such that
\end{enumerate}

\[
h (x) = \int h d \lambda\tag{10}
\]

for every function \(h \in \mathcal{H}\) (which implies in particular that \(\lambda\) has total mass 1, since \(1 \in \mathcal{H}\) ) be the Dirac measure \(\varepsilon_x\) .

It follows that, for each \(h \in \mathcal{H}\) , the function \(z' \mapsto \langle h, z' \rangle\) on \(C\) attains its supremum at at least one extremal point of \(C\) (TVS, II, §7, No.~1, Prop. 1), and this point belongs to \(i_{\mathcal{H}}(X)\) (loc. cit., Cor. of Prop. 2).

\begin{enumerate}
\def\labelenumi{(\roman{enumi})}
\item
  By (9), \(\| i_{\mathcal{H}}(x)\| \leqslant 1\) in the normed space \(\mathcal{H}'\) , in other words \(i_{\mathcal{H}}(\mathrm{X})\) is bounded, and the assertion follows from the fact that \(\mathcal{H}'\) , equipped with the weak topology \(\sigma (\mathcal{H}',\mathcal{H})\) , is quasi-complete (TVS, III, §4, No.~2, Cor. 5 of Th. 1).
\item
  Every positive measure \(\mu\) of mass 1 on \(i_{\mathcal{H}}(X)\) arises, by transport of structure by means of the homeomorphism \(i_{H}\) , from a positive measure \(\lambda\) of mass 1 on X, the Dirac measure \(\varepsilon_{i_{\mathcal{H}(x)}}\) arising from \(\varepsilon_{x}\) . To say that \(\mu\) admits \(i_{\mathcal{H}(x)}\) as barycenter means, by definition, that
\end{enumerate}

\[
\int_ {\mathrm{X}} \langle h, i _ {\mathcal {H}} (z) \rangle d \lambda (z) = \langle h, i _ {\mathcal {H}} (x) \rangle
\]

for every function \(h \in H\) . Taking (9) into account, the assertion (ii) is just the translation of the criterion of No.~2, Cor. of Prop. 3 for \(i_{\mathcal{H}}(x)\) to be an extremal point of C.

We shall say that a point \(x \in \mathbf{X}\) satisfying condition (ii) of Prop. 4 is \(\mathcal{H}\) -extremal; we denote by \(\mathrm{Ch}_{\mathcal{H}}(\mathbf{X})\) (or simply \(\mathrm{Ch}(\mathbf{X})\) ) the set of these points, and by \(\check{\mathrm{S}}_{\mathcal{H}}(\mathbf{X})\) (or simply \(\check{\mathrm{S}}(\mathbf{X})\) ) the closure of \(\mathrm{Ch}_{\mathcal{H}}(\mathbf{X})\) in \(\mathbf{X}\) .

PROPOSITION 5. --- Every function \(h \in \mathcal{H}\) attains its supremum at at least one \(\mathcal{H}\) -extremal point.

Let \(x\) be a point of \(X\) , \(h\) a function in \(\mathcal{H}\) . The relation \(h(z) \leqslant h(x)\) for all \(z \in X\) may be written \(\langle h, i_{\mathcal{H}}(z) \rangle \leqslant \langle h, i_{\mathcal{H}}(x) \rangle\) for all \(z \in X\) , and therefore means that the weakly closed hyperplane of \(\mathcal{H}'\) with equation \(\langle h, t' \rangle = \langle h, i_{\mathcal{H}}(x) \rangle\) is a support hyperplane of \(i_{\mathcal{H}}(X)\) . It is known (TVS, II, §7, No.~1, Cor. of Prop. 1) that such a hyperplane contains at least one extremal point of the closed convex envelope of \(i_{\mathcal{H}}(X)\) , and such a point \(i_{\mathcal{H}}(y)\) is the image of an \(\mathcal{H}\) -extremal point \(y\) by definition; \(h(y)\) is therefore equal to the supremum of \(h\) in \(X\) .

PROPOSITION 6. --- For every point \(x \in X\) , the following properties are equivalent:

\begin{enumerate}
\def\labelenumi{\alph{enumi})}
\item
  \(x\) is \(\mathcal{H}\) -extremal.
\item
  For every open neighborhood U of x in X and every \(\varepsilon > 0\) , there exists a function \(h \geqslant 0\) in H such that \(h(x) \leqslant \varepsilon\) and \(h(y) \geqslant 1\) for every \(y \in X - U\) .
\end{enumerate}

Let \(x\) be any point of \(X\) , \(f\) a function in \(\mathcal{C}(X; \mathbf{R})\) ; it is known (TVS, II, §3, No.~1, Prop. 1) that the infimum of the numbers \(\lambda(f)\) , for all the positive measures on \(X\) such that \(\lambda(h) = h(x)\) for every function \(h \in \mathcal{H}\) , is equal to the supremum of the numbers \(h(x)\) , where \(h\) runs over the set of functions \(h \in \mathcal{H}\) such that \(h \leqslant f\) . Suppose that \(x\) is \(\mathcal{H}\) -extremal; it then follows from Prop. 4, (ii) that for every function \(f \in \mathcal{C}(X; \mathbf{R})\) ,

\[
f (x) = \sup _ {h \in \mathscr {H}, h \leqslant f} h (x).\tag{11}
\]

To show that \(a)\) implies \(b)\) , we take for \(f\) a continuous mapping of X into \([0,1]\) , with support contained in U, such that \(f(x)=1\) ; then, by (11), there exists a function \(h'\in\mathscr{H}\) such that \(h'\leqslant f\) and \(h'(x)\geqslant 1-\varepsilon\) . Since \(1\in\mathscr{H}\) , the function \(h=1-h'\) meets the conditions of \(b)\) .

Conversely, suppose that the condition \(b)\) is verified; this condition implies that \(1 - h \leqslant \varphi_{\mathrm{U}}\) ; for every positive measure \(\lambda\) on \(X\) satisfying the condition (10), we therefore have

\[
\lambda (\mathrm{U}) = \lambda (\varphi_ {\mathrm{U}}) \geqslant \lambda (1 - h) = 1 - h (x) \geqslant 1 - \varepsilon .
\]

Since, by hypothesis, this relation holds for every \(\varepsilon > 0\) and every open neighborhood U of x, it follows that

\[
\lambda (\{x \}) = \inf _ {U} \lambda (U) \geqslant 1 - \varepsilon
\]

for every \(\varepsilon > 0\) , therefore \(\lambda(\{x\}) = 1\) . Since \(\lambda\) is positive and of total mass 1, necessarily \(\lambda = \varepsilon_x\) , which proves that \(x\) is \(\mathcal{H}\) -extremal, by virtue of Prop. 4, (ii).

PROPOSITION 7. --- Let F be a closed subset of X. The following properties are equivalent:

\begin{enumerate}
\def\labelenumi{\alph{enumi})}
\item
  F contains \(\check{\mathrm{S}}_{\mathcal{H}}(\mathrm{X})\)
\item
  For every function \(h \in \mathcal{H}\) , the set \(F\) intersects the set of points of \(X\) where \(h\) attains its supremum.
\item
  For every point \(x \in X\) , there exists a positive measure \(\mu\) of total mass 1 on \(X\) , such that \(\operatorname{Supp}(\mu) \subset F\) and \(h(x) = \int h d\mu\) for every function \(h \in \mathcal{H}\) .
\end{enumerate}

Let \(G = i_{\mathcal{H}}(F)\) . The condition \(a\) ) signifies that \(G\) contains the set of extremal points of \(C\) . The condition \(b\) ) signifies that \(G\) meets the intersection of \(i_{\mathcal{H}}(X)\) with each of the closed support hyperplanes of \(i_{\mathcal{H}}(X)\) . Finally, the condition \(c\) ) signifies that every point of \(i_{\mathcal{H}}(X)\) is the barycenter of a measure with support contained in \(G\) ; by No.~1, Prop. 1, this is also equivalent to saying that the closed convex envelope of \(i_{\mathcal{H}}(X)\) is equal to the closed convex envelope of \(G\) . The equivalence of the conditions \(a\) ), \(b\) ) and \(c\) ) therefore follows from TVS, II, §7, No.~1, Cor. of Prop. 2.

PROPOSITION 8. --- Suppose X is metrizable. Then the set \(\operatorname{Ch}_{\mathcal{H}}(\mathbf{X})\) of H-extremal points of X is the intersection of a countable family of open sets in X, and for every \(x \in X\) , there exists a positive measure \(\mu\) of total mass 1 on X such that

\[
\mu (\mathrm{X} - \mathrm{Ch} _ {\mathcal {H}} (\mathrm{X})) = 0 \text { and } \int h d \mu = h (x)
\]

for every \(h \in \mathcal{H}\) .

This is the translation of Th. 1 of No.~2, by transport of structure by means of the homeomorphism \(x \mapsto i_{\mathcal{H}}(x)\) , as in Prop. 5.

A certain number of results of this No.~may be extended when H is replaced by a set P of functions defined on X, with values in \(R \cup \{+\infty\}\) , that are lower semi-continuous, P being assumed to contain the constants and to satisfy \(P + P \subset P\) (Exer. 2).

*Example. --- Take X to be the unit ball \(\| \mathbf{x} \| \leqslant 1\) in \(\mathbf{R}^3\) , and let \(\mathcal{H}\) be a vector space of continuous functions on X, containing the restrictions to X of the affine linear functions on \(\mathbf{R}^3\) and satisfying the `maximum principle', that is, for every non-constant function \(h \in \mathcal{H}\) , the set of points of X where \(h\) attains its supremum is contained in the sphere \(\mathbf{S}_2\) . It then follows easily from Props. 5 and 7 that \(\mathrm{Ch}_{\mathcal{H}}(\mathrm{X}) = \check{\mathrm{S}}_{\mathcal{H}}(\mathrm{X}) = \mathbf{S}_2\) . An important example of a vector space \(\mathcal{H}\) satisfying the preceding conditions is the set of functions continuous on X and harmonic in the open ball \(\| \mathbf{x} \| < 1\) . For these functions, one proves that the positive measure \(\mu\) of mass 1 such that \(\operatorname{Supp}(\mu) \subset \mathbf{S}_2\) and \(h(\mathbf{x}) = \int h d\mu\) for all \(h \in \mathcal{H}\) is given, if \(\| \mathbf{x} \| < 1\) , by Poisson's formula

\[
d \mu (\mathbf {z}) = \frac {1 - \| \mathbf {z} \| ^ {2}}{\| \mathbf {z} - \mathbf {x} \| ^ {3}} d \sigma (\mathbf {z}),
\]

where \(\sigma\) is the measure on \(\mathbf{S}_2\) invariant under the orthogonal group and such that \(\sigma(\mathbf{S}_2) = 1\) (Ch. VII, §3, Exer. 8).*

